% Pillar-7 appendix: derivations for T1-T3, PCD expectation, micro-jitter
\section{Derivations}
\label{sec:p7-appendix}
Throughout, a group is $G$ conditionally i.i.d.\ binary rewards
$R_1, \dots, R_G \sim \mathrm{Bernoulli}(p)$ for a prompt with latent
success probability $p$; $K = \sum_i R_i \sim \mathrm{Binomial}(G, p)$ and
$\hat p = K/G$.

\subsection{T1: expected ZVF indicator}
\label{sec:p7-appendix-t1}
The within-group sample variance is zero iff all rewards agree, i.e.\
$K = 0$ or $K = G$. These events are disjoint, so
\begin{equation}
  \E[\mathbf{1}\{\text{degenerate}\}]
  \;=\; \Pr(K = 0) + \Pr(K = G)
  \;=\; (1-p)^G + p^G
  \;=\; h_G(p).
\end{equation}
For a prompt population $\mathcal{D}$ with per-prompt success probabilities
$p_x$, linearity of expectation gives
$\E[\mathrm{ZVF}] = \E_{x \sim \mathcal{D}}[h_G(p_x)]$. Note the Jensen gap
used in Section~\ref{sec:p7-t1}: since $h_G$ is convex near the endpoints
and the population mixes prompts of heterogeneous difficulty,
$\E_x[h_G(p_x)] \geq h_G(\E_x[p_x])$ whenever mass sits at the extremes,
which is why measured run-level ZVF exceeds $h_G(\bar p)$ in
Table~\ref{tab:p7-gsweep}. \hfill$\square$

\subsection{T2: monotonicity in $G$}
\label{sec:p7-appendix-t2}
For $p \in (0,1)$, both $p < 1$ and $1-p < 1$, so $p^{G+1} < p^G$ and
$(1-p)^{G+1} < (1-p)^G$ strictly; summing,
$h_{G+1}(p) < h_G(p)$. The decay is exponential with rate set by the larger
term: $h_G(p) = \max(p, 1-p)^G\,[1 + o(1)]$, so
$-\frac{1}{G}\ln h_G(p) \to -\ln \max(p, 1-p)$. Two regimes matter for
control. At $p = \nicefrac12$, $h_G = 2^{1-G}$: each additional rollout
halves the degenerate fraction (an $e$-folding of $h_G$ costs
$1/\ln 2 \approx 1.44$ rollouts). Near the boundary, $p = 1 - \eps$ with
$\eps \ll 1/G$: $h_G(p) = (1-\eps)^G + \eps^G \approx 1 - G\eps$, so the
marginal rollout recovers only $O(\eps)$ of signal---escalation is
exponentially effective in the interior and linearly futile at the edges,
which is Rule~1 of Section~\ref{sec:p7-rules}. \hfill$\square$

\subsection{T3: advantage mass}
\label{sec:p7-appendix-t3}
With unnormalized group-centered advantages $A_i = R_i - \hat p$, a group
with $K$ successes has $K$ advantages equal to $1 - \hat p$ and $G - K$
equal to $-\hat p$. Hence
\begin{equation}
  \frac{1}{G} \sum_{i=1}^{G} |A_i|
  \;=\; \hat p (1 - \hat p) + (1 - \hat p)\, \hat p
  \;=\; 2\, \hat p (1 - \hat p),
\end{equation}
and the summed advantage magnitude is $2 G \hat p(1-\hat p)$, i.e.\
proportional to the within-group contrast. The policy-gradient norm also
carries the score-function factor
$\nabla_\theta \log \pi_\theta(y_i \mid x)$, which is policy- and
sequence-dependent; T3 therefore predicts a monotone association between
$\|\nabla\|$ and $\hat p(1-\hat p)$, not proportionality---this is what E1
tests ($\mathrm{corr} = +0.71$, Section~\ref{sec:p7-e1}). Under
std-normalized advantages (dividing by
$\sqrt{\hat p(1-\hat p)}$), the mean absolute advantage becomes
$2\sqrt{\hat p(1-\hat p)}$: normalization flattens but does not remove the
$p(1-p)$ dependence, and it leaves the degenerate case (undefined $0/0$,
masked in practice) untouched. \hfill$\square$

\subsection{PCD expectation}
\label{sec:p7-appendix-pcd}
The population variance of the group, $\hat p(1-\hat p)$, has expectation
\begin{equation}
  \E[\hat p(1-\hat p)]
  \;=\; \E[\hat p] - \E[\hat p^2]
  \;=\; p - \Bigl(\frac{p(1-p)}{G} + p^2\Bigr)
  \;=\; \Bigl(1 - \frac{1}{G}\Bigr) p (1-p)
  \;=\; \frac{G-1}{G}\, p(1-p),
\end{equation}
using $\mathrm{Var}(\hat p) = p(1-p)/G$. Averaging over prompts yields
Eq.~\eqref{eq:p7-pcd}. Equivalently, $\hat p(1-\hat p)$ equals
$\frac{1}{G^2} \cdot K(G-K)$, i.e.\ (up to the factor $2/G^2$) the number of
discordant correct/incorrect \emph{pairs} in the group---hence
``pairwise-contrast density,'' and hence the connection to the
pairwise-preference reading of group-relative objectives
\citep{wu2025grpo_dpo, rafailov2024direct}: PCD counts, per unit of rollout
budget, exactly the pairs from which a DPO-style contrast could be formed.
\hfill$\square$

\subsection{Micro-jitter proposition}
\label{sec:p7-appendix-jitter}
Let $R'_i = R_i + \eps_i$ with $\eps_i \stackrel{\text{iid}}{\sim}
\mathrm{U}(0, \delta)$, $\delta = 10^{-4}$. (i) \emph{ZVF collapses:} the
jittered rewards are almost surely pairwise distinct (ties have Lebesgue
measure zero), so every group has nonzero sample variance and the tie-based
ZVF indicator is $0$ for all groups. (ii) \emph{PCD is stable:} in
expectation the within-group variance shifts by exactly
$\mathrm{Var}(\eps) = \delta^2/12 \leq 8.4 \times 10^{-10}$; per-group
sample cross terms between $R$ and $\eps$ are $O(\delta)$ in standard
deviation but zero-mean, so they average out over a batch. Empirically
(\texttt{scripts/pcd\_vs\_zvf.py}, $600$ groups): ZVF $0.1583 \to 0.0000$,
PCD $0.153802 \to 0.153802$ (shift below $10^{-6}$). The practical
corollary is stated in Section~\ref{sec:p7-jitter}: any dense sub-reward of
amplitude $\geq$ machine-distinguishable scale performs this transformation
implicitly, so tie-based ZVF must be computed on the binary verifier output
alone, or replaced by PCD. \hfill$\square$

\subsection{Provenance of the empirical tables}
\label{sec:p7-appendix-provenance}
Table~\ref{tab:p7-modelG} and Table~\ref{tab:p7-ushape} aggregate the
368-run audit in W\&B project \texttt{zvf-audit} (790-run corpus; entity
\texttt{arvindcr4-pes-university}); per-model means are computed over all
completed runs of the named model, and the $G$-grid restricts to runs where
only $G$ varies. Table~\ref{tab:p7-gsweep} is
\texttt{experiments/results/groupsize\_zvf\_sweep.tsv} (12 runs: 4 values
of $G$ $\times$ 3 seeds, full per-run records in the adjacent
\texttt{.json}). The binned
shape and jitter numbers are
\texttt{experiments/results/pcd\_vs\_zvf\_shape.tsv} and
\texttt{pcd\_vs\_zvf\_summary.tsv}, regenerable via
\texttt{scripts/pcd\_vs\_zvf.py} from the released per-group reward
tensors. E1 and E3 live in W\&B project \texttt{zvf-colab-experiments}
(run families \texttt{E1\_grad\_signal}, \texttt{E3\_open\_audit}).
